\paragraph{辅助符号链、外侧审计免疫与固定窗口的张量障碍}
在定理 \ref{thm:xi-time-part62d-triple-frobenius-classfunction-tensor-independence} 已将三域 Frobenius 类函数平均压缩为严格张量乘积，而定理 \ref{thm:xi-time-part62e-triple-product-abelianization-c2-four} 与推论 \ref{cor:xi-time-part62ea-A-leyang-threeway-density} 已分别给出 \(C_2^4\) 的二次特征层与一次因子事件的双域联合密度之后，还可继续抽取出下列四条后果与一个结构性猜想。

\begin{corollary}[三条二次符号链 \texorpdfstring{$(\chi_{10},\chi_{\mathrm{LY}},\chi_4)$}{(chi10, chiLY, chi4)} 在 \texorpdfstring{$\{\pm1\}^3$}{pm1^3} 上严格等分布]\label{cor:xi-time-part62eb-three-sign-equidistribution}
沿用定理 \ref{thm:xi-time-part62e-triple-product-abelianization-c2-four} 的记号，并记
\[
\chi_4:=\chi_{4,\mathrm{aux}}=\iota\circ \mathrm{pr}_{C_2}:S_4\times C_2\to\{\pm1\}.
\]
则对任意
\[
(\varepsilon_{10},\varepsilon_{\mathrm{LY}},\varepsilon_4)\in\{\pm1\}^3
\]
都有
\[
\delta\!\bigl(
\chi_{10}(p)=\varepsilon_{10},
\chi_{\mathrm{LY}}(p)=\varepsilon_{\mathrm{LY}},
\chi_4(p)=\varepsilon_4
\bigr)
=
\frac18.
\]
从而任意两条符号被固定以后，第三条仍保持均匀分布。例如
\[
\delta\!\bigl(
\chi_4(p)=+1\mid
\chi_{10}(p)=\varepsilon_{10},
\chi_{\mathrm{LY}}(p)=\varepsilon_{\mathrm{LY}}
\bigr)=\frac12.
\]
\[
\delta\!\bigl(
\chi_4(p)=-1\mid
\chi_{10}(p)=\varepsilon_{10},
\chi_{\mathrm{LY}}(p)=\varepsilon_{\mathrm{LY}}
\bigr)=\frac12.
\]
\end{corollary}
\begin{proof}
取类函数
\[
f(\tau):=\mathbf 1_{\{\operatorname{sgn}(\tau)=\varepsilon_{10}\}},
\qquad
g(\upsilon):=\mathbf 1_{\{\operatorname{sgn}(\upsilon)=\varepsilon_{\mathrm{LY}}\}},
\qquad
h(\eta):=\mathbf 1_{\{\chi_4(\eta)=\varepsilon_4\}}.
\]
由 \(A_{10}\subset S_{10}\) 与 \(A_3\subset S_3\) 都是指数 \(2\) 子群，得
\[
\langle f\rangle_{S_{10}}=\frac12,
\qquad
\langle g\rangle_{S_3}=\frac12.
\]
而在 \(S_4\times C_2\) 中，\(\chi_4\) 只读取 \(C_2\) 坐标，故
\[
\#\{\eta\in S_4\times C_2:\ \chi_4(\eta)=\varepsilon_4\}=24,
\]
于是
\[
\langle h\rangle_{S_4\times C_2}=\frac{24}{48}=\frac12.
\]
代入定理 \ref{thm:xi-time-part62d-triple-frobenius-classfunction-tensor-independence} 即得联合密度为
\[
\frac12\cdot \frac12\cdot \frac12=\frac18.
\]
条件概率公式再由 Bayes 公式立得。证毕。
\end{proof}

\leanverified{paper\_xi\_time\_part62eb\_leyang\_external\_audit\_immunity}
\begin{proposition}[Lee--Yang 分裂型对一切外侧审计事件严格免疫]\label{prop:xi-time-part62eb-leyang-external-audit-immunity}
沿用定理 \ref{thm:xi-time-part62d-triple-frobenius-classfunction-tensor-independence} 的记号，并记
\[
\sigma_{\mathrm{LY}}(p):=\sigma_3(p)\in S_3,
\qquad
\sigma_{4,\mathrm{tot}}(p):=\bigl(\sigma_4(p),\varepsilon_4(p)\bigr)\in S_4\times C_2.
\]
设
\[
\mathcal E\subset S_{10}\times (S_4\times C_2)
\]
为任意非空共轭类并，并记对应的 prime 事件为
\[
E_{\mathcal E}:=
\bigl\{
p:\ (\sigma_{10}(p),\sigma_{4,\mathrm{tot}}(p))\in \mathcal E
\bigr\}.
\]
则对 \(S_3\) 的任意共轭类 \(D\subset S_3\)，都有
\[
\delta\!\bigl(
\sigma_{\mathrm{LY}}(p)\in D\mid E_{\mathcal E}
\bigr)
=
\frac{|D|}{6}.
\]
特别地，\(g\bmod p\) 的三种分裂型
\[
(1)(1)(1),\qquad (1)(2),\qquad (3)
\]
在任意外侧条件 \(E_{\mathcal E}\) 之下仍分别具有条件密度
\[
\frac16,\qquad \frac12,\qquad \frac13.
\]
\end{proposition}
\begin{proof}
记
\[
G_{\mathrm{ext}}:=S_{10}\times (S_4\times C_2).
\]
由定理 \ref{thm:xi-time-part62d-triple-frobenius-classfunction-tensor-independence} 所依赖的三重直积结构，
\[
\Gal(L_{10}L_{\mathrm{LY}}L_4/\QQ)\cong G_{\mathrm{ext}}\times S_3.
\]
因此对任意非空共轭类并 \(\mathcal E\subset G_{\mathrm{ext}}\) 与任意共轭类 \(D\subset S_3\)，Chebotarev 密度定理给出
\[
\delta\!\bigl(E_{\mathcal E}\cap\{\sigma_{\mathrm{LY}}(p)\in D\}\bigr)
=
\frac{|\mathcal E|}{|G_{\mathrm{ext}}|}\cdot \frac{|D|}{6}
=
\delta(E_{\mathcal E})\cdot \frac{|D|}{6}.
\]
对 \(\delta(E_{\mathcal E})>0\) 归一化便得
\[
\delta\!\bigl(
\sigma_{\mathrm{LY}}(p)\in D\mid E_{\mathcal E}
\bigr)
=
\frac{|D|}{6}.
\]
再把 \(S_3\) 的恒等类、换位类与三循环类大小 \(1,3,2\) 代入，即得三种分裂型的条件密度
\[
\frac16,\qquad \frac36=\frac12,\qquad \frac26=\frac13.
\]
证毕。
\end{proof}

\begin{theorem}[十次线性因子事件、辅助二次符号与 Lee--Yang 分裂型的六个显式联合密度]\label{thm:xi-time-part62eb-A-chi4-leyang-six-densities}
沿用推论 \ref{cor:xi-time-part62ea-A-leyang-threeway-density} 的记号，尤其
\[
A:=\{\text{\(P_{10}\bmod p\) 含一次因子}\},
\]
以及
\[
B_{\mathrm{split}},\qquad
B_{\mathrm{single}},\qquad
B_{\mathrm{noroot}},\qquad
B_{\mathrm{root}}=B_{\mathrm{split}}\sqcup B_{\mathrm{single}},
\]
其中 \(B_{\mathrm{split}},B_{\mathrm{single}},B_{\mathrm{noroot}}\) 分别对应 \(g\bmod p\) 的分裂型
\[
(1)(1)(1),\qquad (1)(2),\qquad (3).
\]
再令 \(\chi_4=\chi_{4,\mathrm{aux}}\) 如推论 \ref{cor:xi-time-part62eb-three-sign-equidistribution}。则对每个 \(\varepsilon\in\{\pm1\}\)，有
\[
\delta\!\bigl(A\cap\{\chi_4=\varepsilon\}\cap B_{\mathrm{split}}\bigr)
=
\frac{28319}{537600},
\]
\[
\delta\!\bigl(A\cap\{\chi_4=\varepsilon\}\cap B_{\mathrm{single}}\bigr)
=
\frac{28319}{179200},
\]
\[
\delta\!\bigl(A\cap\{\chi_4=\varepsilon\}\cap B_{\mathrm{noroot}}\bigr)
=
\frac{28319}{268800}.
\]
从而
\[
\delta\!\bigl(A\cap\{\chi_4=\varepsilon\}\cap B_{\mathrm{root}}\bigr)
=
\frac{28319}{134400}.
\]
\end{theorem}
\begin{proof}
由推论 \ref{cor:xi-time-part62ea-A-leyang-threeway-density}，
\[
\delta(A)=\frac{28319}{44800}.
\]
由推论 \ref{cor:xi-time-part62eb-three-sign-equidistribution}，
\[
\delta(\chi_4=\varepsilon)=\frac12.
\]
因此外侧事件
\[
E_\varepsilon:=A\cap\{\chi_4=\varepsilon\}
\]
的密度为
\[
\delta(E_\varepsilon)=\frac{28319}{44800}\cdot \frac12=\frac{28319}{89600}.
\]
而 \(E_\varepsilon\) 只依赖 \(S_{10}\times (S_4\times C_2)\) 因子，故由命题 \ref{prop:xi-time-part62eb-leyang-external-audit-immunity}，在该条件下 Lee--Yang 三种分裂型的条件密度仍分别为
\[
\frac16,\qquad \frac12,\qquad \frac13.
\]
于是
\[
\delta(E_\varepsilon\cap B_{\mathrm{split}})
=
\frac{28319}{89600}\cdot \frac16
=
\frac{28319}{537600},
\]
\[
\delta(E_\varepsilon\cap B_{\mathrm{single}})
=
\frac{28319}{89600}\cdot \frac12
=
\frac{28319}{179200},
\]
\[
\delta(E_\varepsilon\cap B_{\mathrm{noroot}})
=
\frac{28319}{89600}\cdot \frac13
=
\frac{28319}{268800}.
\]
最后，由
\[
B_{\mathrm{root}}=B_{\mathrm{split}}\sqcup B_{\mathrm{single}}
\]
且两者不交，故
\[
\delta(E_\varepsilon\cap B_{\mathrm{root}})
=
\frac{28319}{537600}+\frac{28319}{179200}
=
\frac{28319}{134400}.
\]
证毕。
\end{proof}

\begin{conjecture}[固定有限审计域窗口的张量刚性障碍]\label{conj:xi-time-part62eb-fixed-audit-window-tensor-rigidity}
设
\[
\{L_i/\QQ\}_{i=1}^{r}
\]
为一族固定的 Galois 扩张，满足各 \(L_i/\QQ\) 的伽罗瓦群 \(G_i\) 固定有限，且两两线性互不相交。记
\[
\mathrm{CF}(G_i)
\]
为 \(G_i\) 上的类函数代数。则任何仅由固定多个 Frobenius 因子
\[
\bigl(\mathrm{Frob}_p(L_1/\QQ),\dots,\mathrm{Frob}_p(L_r/\QQ)\bigr)
\]
上的固定类函数、固定二次特征与固定有限布尔组合所生成的 prime-fingerprint 审计体系，在自然密度极限下都只通过有限维交换张量代数
\[
\bigotimes_{i=1}^{r}\mathrm{CF}(G_i)
\]
及其由共轭类并所生成的有限布尔代数起作用。因而，只要域族、群族与线性互不相交模式固定而不随输入规模增长，此类 Chebotarev 指纹审计不足以单独承载 NP-完全型的组合非对称性。若要在 Frobenius 指纹层面逼近 \(P\)–\(NP\) 分离，必须允许审计域链、Galois 群规模或交域失效机制随输入长度增长。
\end{conjecture}

\endinput
